

\vspace*{-5pt}
\section{Introduction}
\vspace*{-5pt}
% Transformers, word embeddings, word2vec, contextual embeddings, architectural biases. Lot of efforts to understand by probing.

Self-attention based Transformer networks \citep{vaswani2017attention} have been at the center of the recent progress on various natural language processing (NLP) tasks, including machine translation \citep{vaswani2017attention}, language modeling \citep{radford2018gpt,radford2019gpt2}, and question answering \citep{devlin2018bert, xlnet2019, roberta2019}. %, just to name a few.
All these tasks involve learning models that map an input sequence of tokens to an output sequence of tokens. Transformers make it feasible to train large models to approximate these sequence-to-sequence functions due to their ability to process the input tokens in a parallel way, as opposed to the sequential %input processing 
nature of RNNs and LSTMs.

%These models are generally trained in two steps: first train a large Transformer network on  a  language modeling objective with a generic corpus, such as Wikipedia, and then fine-tune this network to get a task specific model. These tasks, generally involve learning a mapping from input sequence of tokens, to output sequence of tokens. For example, in machine  translation the model takes in a sequence of words in source language and outputs a sequence of words in the target language.  

A Transformer block consists of two kinds of layers: a self-attention layer and a token-wise feed-forward layer, with skip connections present in both layers. The self-attention layer transforms each input token embedding using a weighted combination of the embeddings of all tokens in the input sequence,
%computes new embeddings of input tokens using a weighted combination of input token embeddings. 
where weights are generated by pairwise dot-products among the input token embeddings. The token-wise feed-forward layer then independently processes each of these modified input token embeddings without any interaction among them. 
Notably, Transformers employ parameter reuse across tokens, as both layers use the same parameters to process each token.
Moreover, Transformers have to rely solely on the pairwise dot-products to capture interaction between the input tokens.
%without explicit hidden states (cf.~RNNs),
%Both layers have parameter reuse, with the same parameters being used to process all the tokens. The Transformer, with no explicit hidden state unlike RNNs, have to rely solely on the pairwise inner products to update embeddings of the input tokens.
%no explicit hidden state

Given the parameter sharing and limited interactions between tokens, it is natural to wonder: what class of sequence-to-sequence functions can the Transformer networks represent? Also, what is the role of the two different kinds of layers? Are both layers needed to obtain the representation power of Transformers? 
In the existing literature, the advantage of Transformers has often been attributed to their capability of computing \textit{contextual} embeddings/mappings of the input, as opposed to fixed word embeddings as in word2vec \citep{mikolov2013distributed}. Is it possible to formalize the notion of contextual mappings? If yes,
%What is a contextual mapping? 
can Transformers actually compute such mappings? Such questions still remain elusive.

In this paper, we provide a mathematical definition of contextual mappings and show that multi-head self-attention layers can indeed compute contextual mappings of the input sequences. We further show that this ability to compute contextual mappings coupled with the value mapping ability of the feed-forward layers makes Transformers universal approximators of any permutation equivariant sequence-to-sequence function. We also improve this result using positional encodings, and show that Transformers can represent any sequence-to-sequence function; i.e., the restriction of permutation equivariance can be removed by positional encodings.

%We note that embeddings computed by Transformers can have other properties....
%However, this has not been proved, and there is no explicit characterization of nature of such contextual embeddings computed by Transformers. 

These results on universal approximation of sequence-to-sequence functions raise a natural question:
%is it possible to replace the self-attention layer with more efficient components that are capable of implementing the contextual mapping; 
is it possible to have a more efficient architecture to compute contextual mappings, consequently, preserving the ability to universally approximate sequence-to-sequence functions? Towards this, we explore other architectures that can implement contextual mappings (to some extent), and experimentally evaluate their performance. In our experiments, we notice that the models that combine these simpler architectures with Transformers have better performance, compared to the standalone Transformers.
%In our experiments we notice that, while these simpler architectures are weaker than Transformers, the models that combine these simpler architectures with Transformers result in improved performance as compared to the standalone Transformers.
We conclude the paper by presenting more discussion and interesting future research directions along these lines.

%Towards the end we discuss other, relatively simpler, architectures that can compute contextual mappings, and experimentally compare their performance with Transformers. We use the BERT setup \citep{devlin2018bert} for our experiments, and present results on a question answering task (SQuAD) \citep{rajpurkar2016squad} and a sentence entailment task (MNLI) \citep{mnli}. In our experiments we notice that, while these simpler architectures are weaker than Transformers, combining them together can lead to stronger models than either of them individually. We present more discussion and interesting future research directions towards the end of Section \ref{sec:discussion}.

\vspace*{-5pt}
\subsection{Summary of our contributions}
\vspace*{-5pt}
%Here, we summarize the key contributions of this paper.
\begin{list}{\textbullet}{\leftmargin=1.1em \itemindent=0em \itemsep=1pt}
    \item We prove that Transformers are universal approximators of continuous and permutation equivariant sequence-to-sequence functions with compact support (Theorem~\ref{thm:univapprox}). We also show that, if Transformers have trainable positional encodings added to the input, then they are universal approximators of continuous sequence-to-sequence functions on a compact domain (Theorem~\ref{thm:anyseq2seq}). 
    \item We formalize the notion of \emph{contextual mappings} and show that the attention layers can compute contextual mappings, where each unique context is mapped to a unique vector (Lemma~\ref{lemma:contextmap}). 
    %Our proof strategy clearly delineates different roles of attention layers and feed-forward layers in universal approximation. Notably, attention layers can implement \emph{contextual mapping}, where each unique context is mapped to unique vectors.
    \item We experimentally evaluate other simpler layers that can compute contextual mappings to some extent, such as bi-linear projections and separable convolutions, and show that substituting some of the self-attention layers with these layers can result in better performance (Section~\ref{sec:discussion}). %We use the BERT setup \citep{devlin2018bert} for our experiments, and present results on a question answering task (SQuAD) \citep{rajpurkar2016squad} and a sentence entailment task (MNLI) \citep{mnli}.
    %We experimentally evaluate the significance of contextual mapping property by comparing architectures designed to compute contextual embeddings, such as bi-linear projections and separable convolutions, with Transformers. %We conclude the paper by presenting more discussion and interesting future research directions towards the end of Section \ref{sec:discussion}.
\end{list}

%%%%%%%%%%%%%%%%%%%%%%%%%%%%%%%%%%%%%%%%%%%%%%%%%%%%%%%%%%%%%%%%%%%%%%%%%%%%%%%
\vspace*{-5pt}
\subsection{Related works \& notation}
\vspace*{-3pt}
%\paragraph{Analysis of attention-based models.}
\noindent{\bf Analysis of attention-based models.}~%Since the advent of attention-based models such as Transformers, there have been numerous attempts trying to understand the role of attention layers in natural language processing models.
Given the popularity of Transformers, there have been numerous works trying to understand the role of attention layers in natural language processing models. One such line of work focuses on probing the output of attention layers to understand the attention mechanism and internal language representation \citep{hewitt2019structural, clark2019does, coenen2019visualizing, vig2019analyzing}. Although these results give valuable insights, a consistent theoretical analysis corroborating these findings is missing. %these results are lack concrete theoretical analysis.

%\paragraph{Universal approximation theorems.}
\noindent{\bf Universal approximation theorems.}~Universal approximation theorems are classical results in neural network theory, dating back many decades \citep{cybenko1989approximation, hornik1991approximation}. These results show that given unbounded width, a one-hidden-layer neural network can approximate arbitrary continuous function with compact support, up to any accuracy. 
Other results focusing on depth appeared more recently \citep{lu2017expressive, hanin2017approximating, lin2018resnet}. In particular, \citet{lu2017expressive, hanin2017approximating} consider fully-connected ReLU networks whose input dimension is $d$, and show that networks with width $d+1$ and unbounded depth are universal approximators of scalar-valued continuous functions. \citet{lin2018resnet} show that a residual network with one hidden neuron per residual block is a universal approximator of scalar-valued functions, given unbounded depth.
Although Transformer networks do have residual connections, due to their heavy parameter sharing, the existing analyses for residual networks do not extend to Transformers.
\citet{sannai2019universal} consider universally approximating permutation invariant/equivariant functions using fully-connected ReLU networks.

%\paragraph{Turing completeness results on Transformers.}
\noindent{\bf Turing completeness results on Transformers.}~Recently, \citet{perez2019turing} have shown that Transformers with infinite precision are Turing complete, which is not the case in finite precision setting~\citep{dehghani2018universal}. We note that Turing completeness deals with computation on formal languages (thus discrete objects), while universal approximation focuses on functions on a continuum. In other words, these are two different concepts; and one does not imply another.

%There has been recent theoretical results studying whether Transformer networks are Turing complete. It has been shown that Transformers with infinite precision are Turing complete \citep{perez2019turing}, but in finite precision settings it is not the case \citep{dehghani2018universal}. Turing completeness is about computation on formal languages (thus discrete objects), while universal approximation is concerned about functions on a continuum. In other words, they are different concepts, and one does not imply another.

% In this paper we prove that Transformers are universal approximators of any permutation equivariant sequence to sequence function. This is quite surprising given their structure with the parameter reuse.  We further show that using positional encodings, Transformers can represent any sequence to sequence function.  Our analysis highlights the key role of self-attention layer in computing the contextual embeddings of input tokens.  


% \begin{itemize}
%     \item Attention models are important. widespread usage and impressive performance of Transformers. 
%     \item Due to its attention structure and shared parameters, expressive power of transformers is not well-understood
%     \item Universal approximation theorems have been classical results in expressive power of neural networks, dating back to 80s.
%     \item Can Transformers represent arbitrary seq-seq functions? - what is the role of attention/FF layers?
%     \item what is the importance of pairwise inner product? What is the role of multihead?
% \end{itemize}


% \setcounter{section}{-1}
% \section{Outline}
% \begin{itemize}
%     \item Intro \& contributions
%     \item Universal approximation for Transformers 
%     \begin{itemize}
%         \item does attention limit model class? permutation equivariant functions: can it represent any permutation equivariant function? 
%         \item contextual mapping: key property for seq-seq. proof that attention satisfies this.
%         \item Efficient memorization: 
%         \item Main Theorem. proof sketch
%         \item Other architectures: additive, bilinear, conv
%         \item Turing completeness.
%     \end{itemize}
%     \item Experiments.    
%     \begin{itemize}
%         \item experiments highlighting that this is a key property.
%     \end{itemize}
% \end{itemize}
